\subsection{Hybrid-Model}
While reviewing the gameplay of the trained \ac{SAC} and \ac{TD3} agents,
we noticed that both agents have their strengths and weaknesses in different situations.
To combine the strengths of both agents, we created a hybrid model.
This approach is inspired by the \ac{MoE} architecture, where multiple experts are combined using a gating network to weight the experts predictions.
But since we've got an continious spatial action space, weighting both experts predictions would not be feasible option.
Thus a \ac{DDQN} \cite{vanhasselt2015deepreinforcementlearningdouble} was implemented as selection gate for sub-agents.
This gate estimates the value of choosing a speciffic sub-agent, given a state:

\begin{gather}
    i^*(s) = \arg\max_{i \in \text{Sub-Agents}} Q(i, s) \\
    a = \pi_{i^*(s)}(s)
\end{gather}



\begin{figure}[ht]
    \centering
    \begin{tikzpicture}[node distance=2cm, auto,>=latex']
        % Define custom colors
        \definecolor{statecolor}{RGB}{173,216,230}    % Light Blue
        \definecolor{diamondcolor}{RGB}{255,182,193}    % Light Pink
        \definecolor{saccolor}{RGB}{255,255,224}          % Light Yellow
        \definecolor{td3color}{RGB}{224,255,255}          % Light Cyan
        \definecolor{actioncolor}{RGB}{255,228,181}       % Moccasin

        % Input: State s
        \node [draw, rectangle, fill=statecolor, minimum width=2.5cm, minimum height=1cm, rounded corners] (state) {State $s$};

        % Decision: argmax (diamond) directly below state
        \node [draw, diamond, fill=diamondcolor, below of=state, node distance=3cm, aspect=2, inner sep=1pt] (argmax)
            {$\arg\max\limits_{i \in \text{Sub-Agents}} Q(i,s)$};

        % Sub-Agents
        \node [draw, rectangle, fill=saccolor, below left of=argmax, node distance=3.5cm, minimum width=2.8cm, minimum height=1cm, rounded corners] (sac)
            {SAC Agent\\$\pi_{\text{SAC}}(s)$};
        \node [draw, rectangle, fill=td3color, below right of=argmax, node distance=3.5cm, minimum width=2.8cm, minimum height=1cm, rounded corners] (td3)
            {TD3 Agent\\$\pi_{\text{TD3}}(s)$};

        % Output: Action a
        \node [draw, rectangle, fill=actioncolor, below of=argmax, node distance=6cm, minimum width=3cm, minimum height=1cm, rounded corners] (action)
            {Action $a$};

        % Arrows
        \draw[->, thick] (state) -- (argmax);
        \draw[->, thick] (argmax) -- node [left] {$i^*(s)=\text{SAC}$} (sac);
        \draw[->, thick] (argmax) -- node [right] {$i^*(s)=\text{TD3}$} (td3);
        \draw[->, thick] (sac) -- (action);
        \draw[->, thick] (td3) -- (action);
    \end{tikzpicture}
    \caption{Hybrid Model with DDQN-based Agent Selection and \ac{SAC} and \ac{TD3} Sub-Agents}
  \end{figure}



The loss function for training the \ac{DDQN} is based on the \ac{TD}-Error,
which is the difference between the predicted and the target Q-value. We investigated three different training strategies for the hybrid model:
\begin{itemize}
    \item Training from scratch:  Training the subagents and the gating network simultaneously from scratch.
    \item Training by fine-tuning the sub-agents: Using sub-agents that are pre-trained independently and further finetuned simulataneously to the gating network.
    \item Training with fixed sub agents: Using subagents that are pre-trained independently and left fixed while training the gating network.
\end{itemize}
